% 分野別類題: 部分積分・置換積分（定積分の定番）
% 計算スピード重点分野のため5問（基本3問＋やや複雑2問）。
% コンパイル: TEXINPUTS="../../../00_common:$TEXINPUTS" latexmk -xelatex problems.tex
\documentclass[a4paper,11pt]{article}
\usepackage{preamble}

\begin{document}

\kamokumei{類題演習 --- 部分積分・置換積分（定積分の定番）}

\shijibun{次の定積分の値を求めよ。導出過程を示せ。（ヒント: 部分積分の公式 $\displaystyle\int_a^b f(x)g'(x)\,dx = \bigl[f(x)g(x)\bigr]_a^b - \int_a^b f'(x)g(x)\,dx$、および適切な置換による変数変換を用いる。）}

\shomon{問1.}{次の定積分の値を求めよ。}
\[
\int_{0}^{\pi/2} x\sin x\,dx
\]

\shomon{問2.}{次の定積分の値を求めよ。}
\[
\int_{0}^{1} x e^{x}\,dx
\]

\shomon{問3.}{次の定積分の値を求めよ。}
\[
\int_{1}^{e} x\ln x\,dx
\]

\shomon{問4.}{次の定積分の値を求めよ。（循環型の部分積分を用いよ。）}
\[
\int_{0}^{\pi} e^{x}\cos x\,dx
\]

\shomon{問5.}{次の定積分の値を求めよ。（適切な置換を用いよ。）}
\[
\int_{0}^{1} x\sqrt{1-x^{2}}\,dx
\]

\end{document}
